% \documentclass{article}
% \usepackage[utf8]{inputenc}
% \usepackage{amsmath}
% \usepackage{graphicx}
% \usepackage{booktabs}
% \usepackage{array}
% \usepackage{geometry}\geometry{a4paper, margin=1in}
% \usepackage{hyperref}
% \usepackage{parskip}

% \title{Architectural Description: Human Benchmark \\ \Large Operational View}
% \author{Generated by AI Software Architect Assistant}
% \date{\today}

% \begin{document}

% \maketitle
% \tableofcontents
% \newpage

\section{Introduction to the Operational View}

The Operational View describes how the Human Benchmark system will be operated,
administered, and supported when running in its production environment. This
view is critical for ensuring that the system functions reliably in day-to-day
operation, addressing aspects such as maintenance, monitoring, and
support.

The Operational Viewpoint focuses on concerns related to the ongoing operation
of the system after deployment. It bridges the gap between system
development and production usage, ensuring that the system can be effectively
managed throughout its lifecycle.

The primary stakeholders for the Operational View include:
\begin{table}[H]
    \centering
    \caption{Primary Stakeholders for the Operational View}
    \label{tab:operational_view_stakeholders}
    \begin{tabular}{|p{0.3\textwidth}|p{0.6\textwidth}|}
        \hline
        \textbf{Stakeholder}  & \textbf{Role in the Operational View}                        \\ \hline
        System Administrators & Responsible for installation, configuration, and maintenance \\ \hline
        Production Engineers  & Handle deployment and operational monitoring                 \\ \hline
        Support Staff         & Provide assistance to users                                  \\ \hline
        Developers            & Need to understand operational considerations                \\ \hline
        Testers               & Verify operational aspects                                   \\ \hline
        Assessors             & Evaluate operational capabilities and risks                  \\ \hline
    \end{tabular}
\end{table}

This view is particularly important for the Human Benchmark system, as it
requires reliable operation to support both single-player and multiplayer
experiences with real-time interaction.

\section{Core Architectural Elements of the Operational View}

\subsection{Installation and Upgrade Process}

The Human Benchmark system employs a containerized deployment approach to
installation and upgrade. This strategy
provides consistency across environments and simplifies the deployment process.

\subsubsection{Installation Components}

The system installation is organized into the following component groups:

\begin{itemize}
    \item \textbf{Frontend Components}: React-based user interface components deployed as container images
    \item \textbf{Backend API Components}: ASP.NET Core services packaged as container images
    \item \textbf{Database Components}: PostgreSQL database schema definitions and migration scripts
    \item \textbf{WebSocket Services}: Real-time communication components for multiplayer features
    \item \textbf{Configuration Files}: Environment-specific configuration for services
\end{itemize}

\subsubsection{Installation Strategy}

The installation process is managed through a CI/CD pipeline that follows these steps:

\begin{figure}[H]
    \centering
    \includegraphics[width=0.8\textwidth]{AutomationAndSections/Sections/7operationalView/Images/Deployment Diagram.drawio.png}
    \caption{Deployment Diagram: Human Benchmark Installation and Upgrade Process Flow}
    \label{fig:deployment_diagram}
\end{figure}

Below is a more detailed version of versioning sequence diagram:

\begin{figure}[H]
    \centering
    \includegraphics[width=0.8\textwidth]{AutomationAndSections/Sections/7operationalView/Images/Versioning.drawio.png}
    \caption{Versioning Sequence Diagram: Detailed Flow of Version Management in Human Benchmark System}
    \label{fig:versioning_sequence_diagram}
\end{figure}

From this diagram, we can see that we have 5 stages:

The five stages of the process are as follows:

\begin{enumerate}
    \item \textbf{Build Application}: The application is built from the source
    code. During this stage, developers have the option to include a specific
    message in their commit to trigger the versioning process, ensuring that the
    build is appropriately tagged and tracked.

    \item \textbf{Test (Integration and Unit)}: Comprehensive testing is
    conducted, including both integration and unit tests, to validate the
    functionality and reliability of the application. This ensures that the
    application meets quality standards before proceeding to the next stages.

    \item \textbf{Dockerize}: The application is packaged into Docker
    containers. This step ensures consistency across environments and simplifies
    deployment by encapsulating all dependencies and configurations within the
    container.

    \item \textbf{Update Database}: Database updates are manually triggered
    using a special commit message. This ensures that schema changes or
    migrations are applied in a controlled manner, minimizing the risk of errors
    during deployment.

    \item \textbf{Deploy Application to MIF OpenNebula}: The final stage
    involves deploying the application to the MIF OpenNebula environment. This
    step ensures that the application is made available in the production
    environment, ready for end-user access.
\end{enumerate}

\subsubsection{Upgrade Process}

The upgrade process is exactly the same as the installation process.
The only difference is that the developer must ensure that the database migration
is applied correctly otherwise the CI/CD pipeline will fail.

\subsubsection{Rollback Strategy}

In case CI/CD pipeline fails during the upgrade process, simpply nothing will be
deployed to the production environment. A red symbol will appear near the commit
to indicate that the CI/CD failed. If CI/CD succeeds but later on issues are
detected, the rollback must be performed manually by the system administrator.

% TODO: DIAGRAM - Installation and Upgrade Process Flow
% % SUGGESTED_TYPE: UML Activity Diagram
% % KEY_ELEMENTS_TO_DEPICT: CI/CD Pipeline, Container Registry, Database Migration, Health Checks, Blue-Green Deployment, Rollback Path
% \begin{figure}[h!]
%     \centering
%     % \includegraphics[width=0.9\textwidth]{images/placeholder_diagram_for_operational_view_installation_process.png}
%     \caption{Illustrative Diagram: Human Benchmark Installation and Upgrade Process Flow (Conceptual Placeholder)}
%     \label{fig:operational_installation_process}
%     \textit{Comment: This diagram should illustrate the complete flow of installation and upgrade processes, including CI/CD pipeline integration, container deployment, database migrations, and rollback mechanisms.}
% \end{figure}

\subsection{Migration Strategy}
We do not have any migration strategy, because we are not migrating from any
previous version of the Human Benchmark system. Moreover we only have one
database instance. As a result, the simple DB update stage in CI/CD pipeline is
sufficient for now.

% \subsubsection{Functional Migration}

% For users transitioning from earlier versions or similar platforms, the Human Benchmark system implements a staged migration approach:

% \begin{itemize}
%     \item \textbf{User Account Migration}: User accounts and authentication credentials are migrated first
%     \item \textbf{Historical Performance Data}: Performance history and achievements are transferred
%     \item \textbf{Custom Game Configurations}: User-created game configurations are preserved
%     \item \textbf{Group Memberships and Relationships}: Social connections and multiplayer groups
% \end{itemize}

% \subsubsection{Data Migration}

% The data migration process addresses the complexity of moving data between different database structures and versions:

% \begin{itemize}
%     \item \textbf{ETL Pipeline}: Custom extraction, transformation, and loading processes for PostgreSQL
%     \item \textbf{Data Quality Assurance}: Validation rules to ensure data integrity during migration
%     \item \textbf{Delta Synchronization}: Mechanism to capture and apply changes made during migration
%     \item \textbf{Performance Optimization}: Bulk loading techniques and parallel processing for large datasets
% \end{itemize}

% For live system migration, the process includes:
% \begin{enumerate}
%     \item Initial bulk data transfer during off-peak hours
%     \item Capturing ongoing changes in a transaction log during migration
%     \item Applying captured changes to synchronize final state
%     \item Verification of data integrity through consistency checks
% \end{enumerate}

% % TODO: DIAGRAM - Data Migration Architecture
% % SUGGESTED_TYPE: Data Flow Diagram
% % KEY_ELEMENTS_TO_DEPICT: Source Database, ETL Pipeline, Data Quality Checks, Transformation Rules, Target Database, Change Capture Mechanism
% \begin{figure}[h!]
%     \centering
%     % \includegraphics[width=0.9\textwidth]{images/placeholder_diagram_for_operational_view_data_migration.png}
%     \caption{Illustrative Diagram: Human Benchmark Data Migration Architecture (Conceptual Placeholder)}
%     \label{fig:operational_data_migration}
%     \textit{Comment: This diagram should illustrate the data migration architecture, showing the flow from source systems through extraction, transformation, quality checks, and loading into the target Human Benchmark database.}
% \end{figure}

\subsection{Operational Monitoring and Control}

Currently we have no monitoring in place, but we plan to use Graphana and
Elasticsearch combined with Prometheus for monitoring and alerting.

\subsubsection{Monitoring Approach}

We have defined a monitoring approach that includes:

\begin{itemize}
    \item \textbf{Infrastructure Monitoring}: Docker commands to check container status
    \item \textbf{Application Performance Monitoring}: Response times, error rates, and throughput metrics
    \item \textbf{Database Monitoring}: Query performance, connection pools, and storage utilization using Aiven Console
\end{itemize}

\subsubsection{Control Capabilities}

System administrators have access to the following control capabilities:

\begin{itemize}
    \item \textbf{Service Management}: Start, stop, and restart individual services using MIF OpenNebula
    \item \textbf{Scaling Controls}: Adjust resource allocation for components under load using MIF OpenNebula
    \item \textbf{Configuration Updates}: Apply configuration changes without redeployment using SSH to the server 
\end{itemize}

% LGTM ABOVE

\subsubsection{Operational Dashboard}

A centralized Graphana or Elasticsearch dashboard provides real-time visibility into system health and performance. Key metrics displayed include:

\begin{itemize}
    \item System health status across all components
    \item Active user count and distribution
    \item Resource utilization and scaling status
    \item Recent alerts(errors)
\end{itemize}

% TODO: DIAGRAM - Operational Monitoring Architecture
% SUGGESTED_TYPE: Component Diagram
% KEY_ELEMENTS_TO_DEPICT: Monitoring Agents, Metrics Collection, Alerting System, Dashboards, Control Interface, Log Aggregation
\begin{figure}[h!]
    \centering
    % \includegraphics[width=0.9\textwidth]{images/placeholder_diagram_for_operational_view_monitoring_architecture.png}
    \caption{Illustrative Diagram: Human Benchmark Operational Monitoring Architecture (Conceptual Placeholder)}
    \label{fig:operational_monitoring}
    \textit{Comment: This diagram should illustrate the monitoring architecture, showing how metrics are collected from system components, aggregated, analyzed, visualized, and how alerts are generated.}
\end{figure}

\subsection{Alerting System}

\subsubsection{Alert Categories}

The Human Benchmark alerting system classifies alerts into the following categories:

\begin{itemize}
    \item \textbf{Technical Alerts}: Infrastructure, application, or database issues
    \item \textbf{Performance Alerts}: Degradation in system performance metrics
    \item \textbf{Capacity Alerts}: Resource utilization approaching defined thresholds
\end{itemize}

\subsubsection{Alert Routing and Prioritization}
Currently we do not have the capibility to route and prioritize alerts using
email or sms, so we will simply display them in the Graphana or Elasticsearch
dashboard.

\subsubsection{Alert Management}

To prevent alert fatigue, the system implements:

\begin{itemize}
    \item Automatic suppression of duplicate alerts
    \item Configurable thresholds based on time of day and user load
    \item Alert history for tracking recurring issues
\end{itemize}

\subsection{Configuration Management}

\subsubsection{Configuration Groups}

The Human Benchmark system organizes configuration into the following groups:

\begin{itemize}
    \item \textbf{Application Configuration}: Core application settings across services (port numbers)
    \item \textbf{Database Configuration}: PostgreSQL settings and connection parameters (names, passwords)
\end{itemize}

\subsubsection{Configuration Value Sets}
We have no plans for this as our user base is +20 users and we do not expect 
any significant load on the system.
% Different configuration profiles are maintained for various operational scenarios:

% \begin{itemize}
%     \item \textbf{Standard Operation}: Balanced settings for normal user load
%     \item \textbf{High Load}: Optimized for peak usage periods
%     \item \textbf{Maintenance Mode}: Settings for system maintenance windows
%     \item \textbf{Reduced Capacity}: Configuration for degraded infrastructure scenarios
% \end{itemize}

\subsubsection{Current Configuration Details}

Currently we have the following configuration details:

\begin{itemize}
    \item \textbf{Backend Configuration}:
          \begin{itemize}
              \item Database connection string for PostgreSQL
              \item HTTPS certificate for secure communication
          \end{itemize}
    \item \textbf{Frontend Configuration}:
          \begin{itemize}
              \item URL to the backend API
          \end{itemize}
\end{itemize}

\subsubsection{Configuration Change Strategy}
In case any configuration changes happen, we will first test the application 
with these changes locally, then we will push these changes to the production 
branch of the repository. After that, we will run the CI/CD pipeline to deploy the configuration changes.
We think that this is sufficient for our simple application.
% Configuration changes follow this process:

% \begin{enumerate}
%     \item Changes are first applied to development environment
%     \item Testing validates configuration in staging environment
%     \item Changes are applied to production through the CI/CD pipeline
%     \item Automated validation tests verify system behavior
%     \item Rollback capability is maintained if issues are detected
% \end{enumerate}

\subsection{Performance Monitoring}

\subsubsection{Key Performance Metrics}

The following metrics are collected to assess system performance:

\begin{itemize}
    \item \textbf{Response Time}: Average and percentile response times for API endpoints
    \item \textbf{Throughput}: Requests processed per second across services
    \item \textbf{Error Rate}: Percentage of requests resulting in errors
    \item \textbf{Database Performance}: Query execution time and connection utilization
    \item \textbf{WebSocket Performance}: Connection reliability and message delivery time
    \item \textbf{Client-Side Metrics}: Page load time and UI responsiveness
\end{itemize}

% LGTM ABOVE 409 page in the software system architecture book

\subsubsection{Performance Data Collection}

Performance data is collected through:

\begin{itemize}
    \item Application instrumentation using OpenTelemetry
    \item Infrastructure metrics from Kubernetes and container runtime
    \item Database performance counters from PostgreSQL
    \item Real user monitoring (RUM) from client applications
    \item Synthetic transactions for continuous availability testing
\end{itemize}

\subsubsection{Performance Analysis}

Collected metrics are used for:

\begin{itemize}
    \item Real-time performance dashboards
    \item Trend analysis to identify gradual degradation
    \item Capacity planning based on usage patterns
    \item Anomaly detection to identify unusual behavior
    \item Performance testing validation
\end{itemize}

% TODO: DIAGRAM - Performance Monitoring Data Flow
% SUGGESTED_TYPE: Data Flow Diagram
% KEY_ELEMENTS_TO_DEPICT: Metrics Sources, Collection Agents, Time Series Database, Analysis Engine, Dashboards, Alerting Integration
\begin{figure}[h!]
    \centering
    % \includegraphics[width=0.9\textwidth]{images/placeholder_diagram_for_operational_view_performance_monitoring.png}
    \caption{Illustrative Diagram: Human Benchmark Performance Monitoring Data Flow (Conceptual Placeholder)}
    \label{fig:operational_performance_monitoring}
    \textit{Comment: This diagram should illustrate how performance metrics flow from source components through collection, storage, analysis, and visualization systems.}
\end{figure}

\subsection{Support Model}

\subsubsection{Support Tiers}

The Human Benchmark system employs a tiered support model:

\begin{itemize}
    \item \textbf{Tier 1}: First-line support handling common user issues and basic troubleshooting
    \item \textbf{Tier 2}: Technical support addressing more complex issues requiring system knowledge
    \item \textbf{Tier 3}: Specialist support involving developers and architecture team
\end{itemize}

\subsubsection{Support Channels}

Users can obtain support through multiple channels:

\begin{itemize}
    \item In-app help center with FAQs and guides
    \item Email support for registered users
    \item Community forums for peer-to-peer assistance
    \item Direct chat support for premium users
\end{itemize}

\subsubsection{Incident Management}

Support incidents follow this workflow:

\begin{enumerate}
    \item Incident reporting through appropriate channel
    \item Classification and prioritization based on impact
    \item Assignment to appropriate support tier
    \item Investigation and resolution
    \item Knowledge base update with solution
    \item Follow-up with affected users
\end{enumerate}

\subsubsection{Support Tools}

The following tools support the incident management process:

\begin{itemize}
    \item Ticketing system for tracking issues
    \item Knowledge base for common problems and solutions
    \item Diagnostic tools for technical staff
    \item Remote assistance capabilities for complex issues
    \item Analytics for identifying recurring issues
\end{itemize}

% TODO: DIAGRAM - Support Process Flow
% SUGGESTED_TYPE: UML Activity Diagram
% KEY_ELEMENTS_TO_DEPICT: Support Channels, Ticketing System, Triage Process, Escalation Paths, Resolution, Knowledge Base Update
\begin{figure}[h!]
    \centering
    % \includegraphics[width=0.9\textwidth]{images/placeholder_diagram_for_operational_view_support_process.png}
    \caption{Illustrative Diagram: Human Benchmark Support Process Flow (Conceptual Placeholder)}
    \label{fig:operational_support_process}
    \textit{Comment: This diagram should illustrate the complete support process from initial user contact through resolution, showing decision points, escalation paths, and knowledge management integration.}
\end{figure}

\subsection{Backup and Restore}

\subsubsection{Backup Strategy}

The Human Benchmark system implements a comprehensive backup strategy:

\begin{itemize}
    \item \textbf{Database Backups}:
          \begin{itemize}
              \item Full database backups daily
              \item Transaction log backups every 15 minutes
              \item Differential backups every 6 hours
          \end{itemize}
    \item \textbf{Configuration Backups}:
          \begin{itemize}
              \item System configuration backed up after each change
              \item Version-controlled configuration repository
          \end{itemize}
    \item \textbf{User-Generated Content}:
          \begin{itemize}
              \item Custom game configurations backed up daily
              \item Incremental backup of user assets
          \end{itemize}
\end{itemize}

\subsubsection{Backup Storage}

Backups are stored with the following considerations:

\begin{itemize}
    \item Primary backups in the same cloud region for fast access
    \item Cross-region replication for disaster recovery
    \item Immutable storage for protection against ransomware
    \item Retention policies aligned with business requirements
\end{itemize}

\subsubsection{Restore Capabilities}

The system supports several restore scenarios:

\begin{itemize}
    \item \textbf{Point-in-Time Recovery}: Using transaction logs to restore to a specific moment
    \item \textbf{Full System Recovery}: Complete restoration of all components
    \item \textbf{Selective Restoration}: Recovery of specific data without full system restore
    \item \textbf{Cross-Environment Restoration}: Ability to restore production data to testing environments (with anonymization)
\end{itemize}

\subsubsection{Data Consistency}

Special attention is given to maintaining consistency across distributed components:

\begin{itemize}
    \item Transaction boundaries are clearly defined for backup consistency
    \item Related data across multiple services is backed up with consistency markers
    \item Restore operations include integrity verification between components
    \item Comprehensive testing of restore procedures on regular intervals
\end{itemize}

\subsection{Operation in Third-Party Environments}

\subsubsection{Cloud Deployment Model}

The Human Benchmark system is designed for deployment in cloud environments:

\begin{itemize}
    \item \textbf{Primary Deployment}: Kubernetes clusters in cloud provider regions
    \item \textbf{Database Services}: Managed PostgreSQL database service
    \item \textbf{Storage Services}: Cloud object storage for assets and backups
    \item \textbf{Content Delivery}: CDN integration for static assets and frontend delivery
\end{itemize}

\subsubsection{Cloud Provider Integration}

The system integrates with cloud provider services:

\begin{itemize}
    \item \textbf{Monitoring Integration}: Exports metrics to cloud monitoring solutions
    \item \textbf{Identity Management}: Integration with cloud IAM services
    \item \textbf{Auto-scaling}: Utilizes cloud provider auto-scaling capabilities
    \item \textbf{Load Balancing}: Leverages cloud load balancing services
\end{itemize}

\subsubsection{Multi-Cloud Considerations}

For resilience and flexibility, the system supports multi-cloud operation:

\begin{itemize}
    \item Abstraction layer for cloud-specific services
    \item Portable deployment configurations using Kubernetes
    \item Data replication strategy across cloud providers
    \item Traffic routing capability between cloud environments
\end{itemize}

% TODO: DIAGRAM - Cloud Deployment Architecture
% SUGGESTED_TYPE: Deployment Diagram
% KEY_ELEMENTS_TO_DEPICT: Cloud Regions, Kubernetes Clusters, Managed Services, CDN, Load Balancers, Data Replication
\begin{figure}[h!]
    \centering
    % \includegraphics[width=0.9\textwidth]{images/placeholder_diagram_for_operational_view_cloud_deployment.png}
    \caption{Illustrative Diagram: Human Benchmark Cloud Deployment Architecture (Conceptual Placeholder)}
    \label{fig:operational_cloud_deployment}
    \textit{Comment: This diagram should illustrate the deployment of the system components across cloud infrastructure, showing relationships between managed services, Kubernetes clusters, and globally distributed resources.}
\end{figure}

\section{Perspective Priorities for the Operational View}

The following table outlines the key architectural perspectives and their relative priorities for the Operational View of the Human Benchmark system:

\begin{table}[H]
    \centering
    \caption{Perspective Priorities for the Operational View of Human Benchmark}
    \label{tab:perspective_priorities_operational}
    \begin{tabular}{p{0.3\textwidth}p{0.15\textwidth}p{0.45\textwidth}}
        % \toprule
        \textbf{Architectural Perspective} & \textbf{Priority} & \textbf{Rationale for Human Benchmark within this View}                                                                                                                                                           \\
        % \midrule
        Accessibility                      & \greencircle Low               & While accessibility is important for the user interface, it has minimal impact on operational aspects beyond ensuring admin tools are accessible.                                                                 \\

        Availability and Resilience        & \redcircle High              & Critical for Human Benchmark as a web-based platform with both single-player and multiplayer functionality, requiring continuous availability and minimal downtime, particularly for real-time multiplayer games. \\

        Development Resource               & \yellowcircle Medium            & Operational procedures must balance sophistication with available resources for maintenance and support, especially considering the CI/CD pipeline and containerized deployment approach.                         \\

        Evolution                          & \redcircle High              & The system must support frequent updates with new games and features, requiring robust installation and upgrade processes without disrupting user experience or compromising data integrity.                      \\

        Internationalization               & \greencircle Low               & Has minimal impact on operational aspects beyond ensuring error messages and admin interfaces support multiple languages if deployed internationally.                                                             \\

        Location                           & \yellowcircle Medium            & Relevant for cloud deployment across multiple regions to support global user base with low latency, particularly important for real-time multiplayer functionality.                                               \\

        Performance and Scalability        & \redcircle High              & Critical for Human Benchmark as reaction-time and reflex tests require precise timing and responsiveness, with monitoring and configuration management playing key roles in maintaining performance.              \\

        Regulation                         & \yellowcircle Medium            & Data protection regulations (e.g., GDPR) impact backup, restore, and data migration strategies, particularly regarding user data handling and storage.                                                            \\

        Security                           & \redcircle High              & Essential for protecting user accounts, game integrity, and preventing manipulation of scores and leaderboards, with monitoring and alerting for security incidents being vital.                                  \\

        Usability                          & \yellowcircle Medium            & Administrative interfaces and operational tools must be usable for system administrators and support staff, though not as critical as user-facing components.                                                     \\
        % \bottomrule
    \end{tabular}
\end{table}

\section{Impact of Key Perspectives on the Operational View}

\subsection{Impact of Availability and Resilience Perspective}

Applying the Availability and Resilience perspective to the Operational View of the Human Benchmark system has several significant impacts:

\begin{itemize}
    \item \textbf{Enhanced Monitoring}: The monitoring system has been expanded to include detailed health checks for all components, with particular emphasis on WebSocket connections that support real-time multiplayer functionality. This ensures that any potential issues affecting system availability are detected quickly.

    \item \textbf{Robust Backup Strategy}: The backup and restore capabilities have been designed with higher frequency and multiple storage locations to minimize data loss in failure scenarios. The need for point-in-time recovery capabilities emerged from applying this perspective to ensure that user progress and game history can be reliably restored.

    \item \textbf{Cloud Deployment Architecture}: The third-party environment operation model has been influenced to incorporate redundancy across cloud regions and providers. This includes active-active deployment strategies rather than simpler active-passive approaches, allowing for immediate failover without service interruption.

    \item \textbf{Installation and Upgrade Process}: The blue-green deployment strategy was specifically chosen to minimize downtime during upgrades, with automated rollback capabilities to quickly recover from failed deployments without user impact.
\end{itemize}

\subsection{Impact of Performance and Scalability Perspective}

The Performance and Scalability perspective has significantly shaped the Operational View in the following ways:

\begin{itemize}
    \item \textbf{Comprehensive Performance Monitoring}: The monitoring system has been designed to track fine-grained metrics specifically related to reaction time measurement accuracy, which is essential for the core functionality of Human Benchmark. This includes client-side metrics to ensure timing precision across diverse user devices.

    \item \textbf{Configuration Management}: Multiple configuration value sets have been defined to support different load scenarios, allowing the system to adapt its behavior during peak usage periods without manual intervention. This is particularly important for handling viral traffic spikes that online games often experience.

    \item \textbf{Auto-scaling Design}: The cloud deployment model incorporates horizontal scaling capabilities with predefined thresholds based on key performance indicators, ensuring that resources are automatically provisioned to maintain responsiveness under varying load.

    \item \textbf{Database Performance Optimization}: The operational procedures include specific monitoring and maintenance routines for database optimization, such as index maintenance and query performance analysis, to ensure that leaderboard and historical data queries remain efficient as the dataset grows.
\end{itemize}

\subsection{Impact of Security Perspective}

Applying the Security perspective to the Operational View has resulted in these key impacts:

\begin{itemize}
    \item \textbf{Enhanced Alerting System}: The alerting system has been expanded to include security-specific alerts for unusual patterns of activity, potential cheating attempts, and unauthorized access attempts. These alerts are prioritized to ensure immediate response to security threats.

    \item \textbf{Backup Protection}: The backup strategy has been modified to include immutable storage and encryption to protect against ransomware and data theft, ensuring that backups cannot be compromised even if the primary system is breached.

    \item \textbf{Secure Configuration Management}: The configuration change strategy has been enhanced with approval workflows and audit trails to prevent unauthorized changes that could compromise security. Sensitive configuration values are treated differently with enhanced access controls.

    \item \textbf{Operational Monitoring}: The monitoring approach now includes security-focused telemetry, such as authentication patterns, session anomalies, and unusual API call sequences that might indicate attempts to manipulate game results or leaderboards.
\end{itemize}

\section{Importance and Justification of the Operational View for Human Benchmark}

The Operational View is critically important for the Human Benchmark system for several compelling reasons:

\subsection{System Nature and Goals}

Human Benchmark is a web-based platform that measures users' cognitive and motor skills through precision timing tests and competitive games. The nature of this system makes operational considerations particularly vital:

\begin{itemize}
    \item \textbf{Precision Requirements}: As a system measuring reaction times in milliseconds, operational reliability and performance consistency are fundamental to its core functionality.

    \item \textbf{Dual-Mode Operation}: Supporting both single-player and multiplayer experiences creates complex operational requirements for availability and resilience.

    \item \textbf{Competitive Element}: With leaderboards and user rankings, the integrity of the operational environment directly affects the trustworthiness of the system's primary purpose.

    \item \textbf{Continuous Service}: As a web application with global accessibility, the expectation is for near-continuous availability with minimal downtime.
\end{itemize}

\subsection{Stakeholder Concerns}

The Operational View directly addresses critical concerns of multiple key stakeholders:

\begin{itemize}
    \item \textbf{System Administrators}: Need clear processes for installation, configuration, monitoring, and maintenance to ensure reliable system operation.

    \item \textbf{Production Engineers}: Require well-defined deployment procedures and monitoring capabilities to maintain system health and performance.

    \item \textbf{Support Staff}: Need effective tools and processes to address user issues and system problems quickly.

    \item \textbf{End Users}: Expect consistent performance, data integrity, and availability, particularly for competitive features and historical performance tracking.

    \item \textbf{Business Stakeholders}: Require reliable uptime and scalability to maintain user engagement and satisfaction.
\end{itemize}

\subsection{Quality Attributes Support}

The Operational View is instrumental in achieving several critical quality attributes for Human Benchmark:

\begin{itemize}
    \item \textbf{Performance}: Through comprehensive monitoring, configuration management, and tuning capabilities that ensure the system delivers precise timing measurements.

    \item \textbf{Availability}: Via robust installation, backup/restore, and monitoring processes that minimize downtime and maintain service continuity.

    \item \textbf{Reliability}: Through alerting systems, support processes, and operational procedures that quickly address issues before they impact users.

    \item \textbf{Scalability}: By defining how the system adapts to changing load conditions, particularly important for handling viral growth or promotional events.

    \item \textbf{Security}: Through operational monitoring, alerting, and response procedures that protect user data and system integrity.
\end{itemize}

\subsection{Architectural Principles Alignment}

This view's importance aligns with key architectural principles from the literature:

\begin{itemize}
    \item \textbf{Stakeholder Focus}: The Operational View directly addresses the needs of operational stakeholders who are often overlooked in purely functional or developmental views.

    \item \textbf{Quality Property Support}: It provides concrete mechanisms for achieving critical quality attributes that cannot be addressed through code alone.

    \item \textbf{Practical Implementation}: It bridges the gap between theoretical design and practical reality, addressing the "last mile" of delivering value to users.

    \item \textbf{System Evolution}: It provides the foundation for smooth system evolution through well-defined upgrade and migration processes.
\end{itemize}

In summary, the Operational View is critically important for the Human Benchmark system as it directly impacts the system's ability to deliver on its core promise of accurate performance measurement and competitive gaming. Without robust operational capabilities, even the most elegant functional design would fail to deliver reliable service to users.

\section{Cross-View Consistency Considerations}

Maintaining consistency between the Operational View and other architectural views is essential for the Human Benchmark system to function as intended. This view does not exist in isolation but must align with and complement other views in the comprehensive Architectural Description.

Key cross-view consistency requirements include:

\begin{itemize}
    \item \textbf{Functional View Alignment}: The operational monitoring and control capabilities must align with the functional elements defined in the Functional View. Each functional component should have corresponding operational support for monitoring, configuration, and management.

    \item \textbf{Deployment View Integration}: The installation model and cloud deployment architecture must be consistent with the physical deployment topologies defined in the Deployment View, including server configurations, network connections, and hardware requirements.

    \item \textbf{Information View Consistency}: The backup and restore strategies must account for all data entities and relationships defined in the Information View, ensuring complete data protection and consistency during recovery operations.

    \item \textbf{Development View Coordination}: The operational tools, monitoring interfaces, and support systems must be built according to the module organization and coding standards defined in the Development View.

    \item \textbf{Concurrency View Alignment}: The operational monitoring of real-time multiplayer features must align with the concurrency models defined for simultaneous user interactions and game state management.
\end{itemize}

Changes to any view may necessitate updates to the Operational View to maintain this consistency. For example:

\begin{itemize}
    \item Adding new functional components requires corresponding updates to monitoring, configuration, and support models.

    \item Changes in deployment topology require updates to installation procedures and performance monitoring configurations.

    \item Modifications to data models require adjustments to backup and migration strategies.
\end{itemize}

To facilitate this consistency, we recommend:

\begin{itemize}
    \item Regular cross-view reviews during architectural updates
    \item Explicit traceability between related elements across views
    \item Collaborative design sessions involving stakeholders from different operational domains
    \item Validation of cross-view consistency through integrated testing
\end{itemize}

This document is one component of a comprehensive Architectural Description for the Human Benchmark system, and should be read and maintained in conjunction with the other architectural views to ensure a complete understanding of the system architecture.

% \end{document}